% id: 83-09
% book: 베이직쎈 확률과 통계 · 05 확률변수와 확률분포
% section: 개념 24 이산확률변수와 확률질량함수
% figure: none
% answer: 
% answer_source: 
% variant_level: 0 (원본 전사)
% 이 파일은 build.mjs 가 items.json 에서 만든다. 고칠 때는 items.json 을 고친다.
\smash{\raisebox{1.5mm}{\dmkichul{베이직쎈 확률과 통계 83쪽 09번}}}
한 개의 동전을 세 번 던지는 시행에서 뒷면이 나오는 횟수를 확률변수 $X$라 하자. 다음을 구하시오.
\dmsubprob{1} $X$의 확률질량함수\\[3pt] $\mathrm{P}(X=x)={}_{\text{\scriptsize\setlength{\fboxsep}{1.5pt}\framebox[1.5em]{\rule{0pt}{1.6ex}}}}\mathrm{C}_x\left(\blank\right)^{x}\left(\dfrac{1}{2}\right)^{\text{\scriptsize\setlength{\fboxsep}{1.5pt}\framebox[1.5em]{\rule{0pt}{1.6ex}}}}{}^{-x}$\\[5pt] $\hphantom{\mathrm{P}(X=x)}={}_{\text{\scriptsize\setlength{\fboxsep}{1.5pt}\framebox[1.5em]{\rule{0pt}{1.6ex}}}}\mathrm{C}_x\left(\blank\right)^{3}$ $(x=0{,}\mkern3mu\mathopen{}\blank{,}\mkern3mu\mathopen{}2{,}\mkern3mu\mathopen{}\blank)$
\dmsubprob{2} $X$의 확률분포를 나타낸 표\par\smallskip\centerline{{\setlength{\tabcolsep}{7pt}\renewcommand{\arraystretch}{1.9}\begin{tabular}{|c|c|c|c|c|c|}\hline $X$ & \makebox[2.6em]{} & \makebox[2.6em]{} & \makebox[2.6em]{} & \makebox[2.6em]{} & 합계 \\ \hline $\mathrm{P}(X=x)$ & & & & & \\ \hline\end{tabular}}}
